\chapter{Комплекс функцийн хязгаар ба тасралтгүй чанар}

\section{Комплекс хувьсагчийн функц}

\begin{dfn}
$D\subset\mathbb{C}$ олонлог дээр $f$ дүрэм өгөгдөж, $z\in D$ бүрт комплекс тоо $w$-г
$w=f(z)$ гэж харгалзуулж байвал $f$-г \emph{комплекс хувьсагчийн функц} гэнэ:
\[
f: D\to \mathbb{C},\qquad z\mapsto w=f(z).
\]
$D$-г \emph{тодорхойлогдох муж} гэнэ.
\end{dfn}

Хэрэв $z=x+iy$, $w=u+iv$ ($x,y,u,v\in\mathbb{R}$) гэвэл
\[
f(z)=u(x,y)+iv(x,y),
\]
өөрөөр хэлбэл комплекс функц нь хоёр хувьсагчийн бодит утгат $u(x,y)$, $v(x,y)$ функцүүдээр тодорхойлогдоно:
\[
u(x,y)=\Re f(z),\qquad v(x,y)=\Im f(z).
\]

\begin{exm}
$w=z^2$ функцийн бодит ба хуурмаг хэсгийг ол.
\end{exm}

\begin{sol}
$z=x+iy$ гэвэл
\[
z^2=(x+iy)^2=x^2-y^2+2xy\,i.
\]
Иймээс
\[
u(x,y)=x^2-y^2,\qquad v(x,y)=2xy.
\]
\end{sol}

%---------------------------------------------------

\section{Функцийн хязгаар}

\begin{dfn}
$z_0\in D$ байг. Хэрэв $A\in\mathbb{C}$ тоо олдож
\[
\forall \varepsilon>0\ \exists \delta>0:\ 
0<|z-z_0|<\delta \Rightarrow |f(z)-A|<\varepsilon
\]
нөхцөл биелбэл $A$-г $f(z)$-ийн $z\to z_0$ үед авах \emph{хязгаар} гэнэ.
Тэмдэглэгээ:
\[
\lim_{z\to z_0} f(z)=A.
\]
\end{dfn}

\begin{thm}[Хязгаарын бодит ба хуурмаг хэсэг]
$f(z)=u(x,y)+iv(x,y)$, $z_0=x_0+iy_0$, $A=\alpha+i\beta$ гэе.
Тэгвэл
\[
\lim_{z\to z_0} f(z)=A
\]
байх нь
\[
\lim_{(x,y)\to(x_0,y_0)}u(x,y)=\alpha,\qquad
\lim_{(x,y)\to(x_0,y_0)}v(x,y)=\beta
\]
байхтай эквивалент.
\end{thm}

\begin{proof}
$(\Rightarrow)$ гэж үзье. 
\[
\lim_{z\to z_0} f(z)=A
\]
гэж авъя. Тэгвэл $\forall \varepsilon>0$ тоо бүрт $\exists \delta>0$ олдож
\[
0<|z-z_0|<\delta \ \Rightarrow\ |f(z)-A|<\varepsilon
\]
биелнэ.

Энд
\[
f(z)-A=(u(x,y)-\alpha)+i(v(x,y)-\beta)
\]
учир
\[
|u(x,y)-\alpha|\le |f(z)-A|,\qquad |v(x,y)-\beta|\le |f(z)-A|
\]
гэж үнэлж болно. Иймээс дээрх нөхцлөөс
\[
0<|z-z_0|<\delta \ \Rightarrow\ |u(x,y)-\alpha|<\varepsilon
\]
мөн
\[
0<|z-z_0|<\delta \ \Rightarrow\ |v(x,y)-\beta|<\varepsilon
\]
гарна. Харин $|z-z_0|=\sqrt{(x-x_0)^2+(y-y_0)^2}$ тул энэ нь яг
\[
\lim_{(x,y)\to(x_0,y_0)}u(x,y)=\alpha,\qquad
\lim_{(x,y)\to(x_0,y_0)}v(x,y)=\beta
\]
гэдгийг илэрхийлнэ.

$(\Leftarrow)$ Одоо эсрэгээр нь
\[
\lim_{(x,y)\to(x_0,y_0)}u(x,y)=\alpha,\qquad
\lim_{(x,y)\to(x_0,y_0)}v(x,y)=\beta
\]
гэж үзье. Тэгвэл $\forall \varepsilon>0$ тоо бүрт $\exists \delta_1>0$, $\exists \delta_2>0$ олдож
\[
0<\sqrt{(x-x_0)^2+(y-y_0)^2}<\delta_1 \Rightarrow |u(x,y)-\alpha|<\frac{\varepsilon}{\sqrt2},
\]
\[
0<\sqrt{(x-x_0)^2+(y-y_0)^2}<\delta_2 \Rightarrow |v(x,y)-\beta|<\frac{\varepsilon}{\sqrt2}
\]
биелнэ. $\delta=\min\{\delta_1,\delta_2\}$ гэж авбал
\[
0<|z-z_0|<\delta \Rightarrow 
|u(x,y)-\alpha|<\frac{\varepsilon}{\sqrt2},\quad
|v(x,y)-\beta|<\frac{\varepsilon}{\sqrt2}.
\]

Тэгвэл
\[
|f(z)-A|^2
=
|(u-\alpha)+i(v-\beta)|^2
=
(u-\alpha)^2+(v-\beta)^2
<
\frac{\varepsilon^2}{2}+\frac{\varepsilon^2}{2}
=\varepsilon^2.
\]
Иймээс
\[
0<|z-z_0|<\delta \Rightarrow |f(z)-A|<\varepsilon.
\]
Энэ нь
\[
\lim_{z\to z_0} f(z)=A
\]
гэдгийг баталж байна.
\end{proof}


\begin{sen}
Хэрэв $f(z)$-ийн хязгаар орших бол $z_0$ руу ямар ч замаар ойртсон
утга ижил байна.
\end{sen}

\begin{exm}
\[
\lim_{z\to 0} \frac{z}{\overline{z}}
\]
орших уу?
\end{exm}

\begin{sol}
$z=re^{i\theta}$ гэж бичвэл
\[
\frac{z}{\overline{z}}=e^{2i\theta}.
\]
$z\to 0$ гэдэг нь $r\to 0$ гэсэн үг. Гэтэл $\dfrac{z}{\overline{z}}=e^{2i\theta}$
нь $\theta$-оос хамаарна. Иймээс хязгаар $r\to 0$ үед $\theta$ хамаарсан тул хязгаар оршихгүй.
\end{sol}

\begin{exm}
\[
\lim_{(x,y)\to(0,0)} \frac{x^2}{x^2+y^2}
\]
орших уу?
\end{exm}

\begin{sol}
$y=0$ үед утга $1$,
$x=0$ үед утга $0$.
Хоёр өөр утга гарч байна.
Иймээс хязгаар оршихгүй.
\end{sol}

\begin{sen}[Шахах зарчим]
Хэрэв
\[
|f(z)|\le g(z)
\]
мөн
\[
\lim_{z\to z_0} g(z)=0
\]
бол
\[
\lim_{z\to z_0} f(z)=0.
\]
\end{sen}

\begin{exm}
\[
\lim_{z\to 0} z^2 \sin \frac{1}{z}
\]
утгыг ол.
\end{exm}

\begin{sol}
\[
|z^2\sin \tfrac{1}{z}| \le |z|^2.
\]
\[
\lim_{z\to 0} |z|^2=0.
\]
Иймээс
\[
\lim_{z\to 0} z^2 \sin \frac{1}{z}=0.
\]
\end{sol}
\subsection*{Хязгаар $z\to\infty$ үед}

\begin{dfn}
Хэрэв
\[
\forall \varepsilon>0\ \exists R>0:\ |z|>R \Rightarrow |f(z)-A|<\varepsilon
\]
бол
\[
\lim_{z\to\infty} f(z)=A.
\]
\end{dfn}

\begin{exm}
\[
\lim_{z\to\infty} \frac{1}{z}
\]
утгыг ол.
\end{exm}

\begin{sol}
\[
\left|\frac{1}{z}\right|=\frac{1}{|z|}\to 0.
\]
Иймээс
\[
\lim_{z\to\infty} \frac{1}{z}=0.
\]
\end{sol}
\begin{thm}[Хязгаарын арифметик үйлдлүүд]
Хэрэв
\[
\lim_{z\to z_0} f(z)=A,\qquad
\lim_{z\to z_0} g(z)=B
\]
бол дараах тэнцэтгэлүүд биелнэ:
\begin{enumerate}
    \item ямар ч тогтмол $c\in\mathbb{C}$-ийн хувьд.
$\lim\limits_{z\to z_0} c\,f(z)=cA$

\item $\lim\limits_{z\to z_0} (f(z)+g(z))=A+B$
\item $\lim\limits_{z\to z_0} (f(z)-g(z))=A-B$
\item$\lim\limits_{z\to z_0} (f(z)g(z))=AB$
\item $\lim\limits_{z\to z_0} \dfrac{f(z)}{g(z)}=\dfrac{A}{B} \qquad B\neq 0.$
\end{enumerate}
\end{thm}
\begin{exm}
Дараах хязгаарыг батал:
\[
\lim_{z\to z_0}\sin z=\sin z_0.
\]
\end{exm}

\begin{sol}
$z=x+iy$, $z_0=x_0+iy_0$ гэж бичье.
Дараах томьёог ашиглая:
\[
\sin(x+iy)=\sin x\,\ch y+i\cos x\,\sh y.
\]

Тэгвэл
\[
\sin z-\sin z_0
=
(\sin x\,\ch y-\sin x_0\,\ch y_0)
+i(\cos x\,\sh y-\cos x_0\,\sh y_0).
\]

Бодит хэсгийг задлавал
\[
\sin x\,\ch y-\sin x_0\,\ch y_0
=
(\sin x-\sin x_0)\ch y
+
\sin x_0(\ch y-\ch y_0).
\]

$x\to x_0$, $y\to y_0$ үед
\[
\sin x-\sin x_0\to 0,\qquad
\ch y-\ch y_0\to 0.
\]
Иймээс бодит хэсэг 0 рүү тэмүүлнэ.

Хуурмаг хэсгийг мөн адил задлахад
\[
\cos x\,\sh y-\cos x_0\,\sh y_0
=
(\cos x-\cos x_0)\sh y
+
\cos x_0(\sh y-\sh y_0),
\]
эндээс мөн 0 гарна.

Иймээс
\[
\lim_{z\to z_0}\sin z=\sin z_0.
\]
\end{sol}
\begin{thm}[Давхар функцийн хязгаар]
Хэрэв
\[
\lim_{z\to z_0} f(z)=w_0,
\qquad
\lim_{w\to w_0} g(w)=L,
\]
бол
\[
\lim_{z\to z_0} g(f(z)) = L.
\]
\end{thm}
\begin{exm}
\[
\lim_{z\to 1} \sin(z^2)
\]
утгыг ол.
\end{exm}

\begin{sol}
\[
\lim_{z\to 1} z^2 = 1.
\]
Синус функц бүх цэгт тасралтгүй тул
\[
\lim_{z\to 1} \sin(z^2)=\sin(1).
\]
\end{sol}

%---------------------------------------------------

\section{Тасралтгүй чанар}

\begin{dfn}
Хэрэв
\[
\lim_{z\to z_0} f(z)=f(z_0)
\]
бол $f$-г $z_0$ дээр тасралтгүй гэнэ.
\end{dfn}

\begin{thm}
$f(z)=u(x,y)+iv(x,y)$ бол $f$ нь $z_0$ дээр тасралтгүй байх нь
\[
u(x,y)\ \text{ба}\ v(x,y)
\]
нь $(x_0,y_0)$ дээр тасралтгүй байхтай эквивалент.
\end{thm}

\begin{thm}
Хэрэв $f$ ба $g$ нь $z_0$ дээр тасралтгүй бол
\[
f+g,\quad fg,\quad \frac{f}{g}\ (g(z_0)\neq 0)
\]
нь мөн $z_0$ дээр тасралтгүй.
\end{thm}

\begin{thm}

$P(z)$ олон гишүүнт нь $\mathbb{C}$ дээр бүх цэгт тасралтгүй.
Рациональ функц
\[
R(z)=\frac{P(z)}{Q(z)}
\]
нь
\[
Q(z)\neq 0
\]
байх бүх цэгт тасралтгүй.
\end{thm}
\begin{thm}
Комплекс экспонент функц
\[
f(z)=e^z
\]
нь $\mathbb{C}$ хавтгай дээр бүх цэгт тасралтгүй.
\end{thm}

\begin{proof}
$z=x+iy$ гэж бичвэл
\[
e^z=e^{x+iy}=e^x(\cos y+i\sin y).
\]

Иймээс
\[
u(x,y)=e^x\cos y,\qquad v(x,y)=e^x\sin y.
\]

Бодит функцууд $e^x,\cos y,\sin y$ нь $\mathbb{R}$ дээр тасралтгүй.
Үржвэр нь мөн тасралтгүй тул $u(x,y)$ ба $v(x,y)$ нь
$(x_0,y_0)$ дээр тасралтгүй.

Тэгвэл “$u,v$ тасралтгүй $\Rightarrow$ $f$ тасралтгүй” шалгуураар
$e^z$ нь бүх $z_0\in\mathbb{C}$ дээр тасралтгүй.
\end{proof}

\begin{thm}
Логарифм
\[
\ln z=\ln|z|+i\arg z
\]
нь
\[
\Omega=\mathbb{C}\setminus\{x\le 0,\ y=0\}
\]
муж дээр тасралтгүй.
\end{thm}

\begin{proof}
$z=x+iy$ гэж бичье.

Тэгвэл
\[
\ln z=\frac12\ln(x^2+y^2)
+i\,\arg(z).
\]

1) $\ln(x^2+y^2)$ нь $(0,0)$-оос бусад бүх цэг дээр тасралтгүй.

2) $\Arg(z)$ нь сөрөг бодит тэнхлэгийг хассан муж дээр тасралтгүй.

Иймээс бодит ба хуурмаг хэсгүүд нь $\Omega$ муж дээр тасралтгүй.

“$u,v$ тасралтгүй $\Rightarrow$ $f$ тасралтгүй” шалгуураар
$\Ln z$ нь $\Omega$ муж дээр тасралтгүй.
\end{proof}
\section{Олон утгат функц}

\begin{dfn}
Хэрэв комплекс тоо $z$ бүрт нэгээс илүү $w$ утга харгалзаж байвал
түүнийг \emph{олон утгат комплекс функц} гэнэ.
\end{dfn}
\begin{exm}
Комплекс тооны аргумент.
\end{exm}

Комплекс тоо $z\neq 0$-г
\[
z=re^{i\varphi}
\]
гэж бичвэл
\[
\arg z=\varphi+2k\pi,\qquad k\in\mathbb{Z}.
\]

Иймээс $\arg z$ нь төгсгөлгүй олон утгатай функц.
\begin{exm}
Комплекс логарифм.
\end{exm}

$z\neq 0$ үед
\[
\Ln z=\ln|z|+i(\arg z+2k\pi),\qquad k\in\mathbb{Z}.
\]

Иймээс $\Ln z$ нь олон утгат функц.
\begin{exm}
Комплекс язгуур.
\end{exm}

$z=re^{i\varphi}$ гэж бичвэл
\[
\sqrt[n]{z}
=
r^{1/n}
e^{\,i\frac{\varphi+2k\pi}{n}},
\qquad k=0,1,\dots,n-1.
\]

Иймээс $n$-зэрэгт язгуур нь $n$ утгатай функц.
\begin{dfn}
Олон утгат функцийн утгуудын дундаас нэгийг нь
сонгон авснаар үүсэх нэг утгат функцийг
\emph{салбар} гэнэ.
\end{dfn}
\begin{dfn}
Хэрэв
\[
-\pi<\Arg z\le \pi
\]
гэж сонговол үүсэх логарифмыг
\emph{логарифмын гол утга} гэнэ:
\[
\ln z=\ln|z|+i\Arg z.
\]
\end{dfn}
\section{Семинар}
\begin{enumerate}

\item $f(z)=z^3$ функцийн бодит ба хуурмаг хэсгийг ол.

\item $f(z)=e^z$ функцийн бодит ба хуурмаг хэсгийг ол.

\item $\displaystyle \lim_{z\to 2} (3z^2-1)$ утгыг ол.

\item $\displaystyle \lim_{z\to 1+i} z^2$ утгыг ол.

\item $\displaystyle \lim_{z\to 0} \frac{z}{z+1}$ утгыг ол.

\item $\displaystyle \lim_{z\to 0} \frac{z^2+2z}{z}$ утгыг ол.

\item $\displaystyle \lim_{z\to z_0} \cos z$ утгыг ол.

\item $\displaystyle \lim_{z\to z_0} e^z$ утгыг ол.

\item $\displaystyle \lim_{z\to \infty} \frac{1}{z}$ утгыг ол.

\item $\displaystyle \lim_{z\to \infty} \frac{2z+1}{z}$ утгыг ол.

\item Дараах хязгаар орших эсэхийг шалга:
\[
\lim_{z\to 0} \frac{|z|}{z}.
\]

\item Дараах хязгаар орших эсэхийг шалга:
\[
\lim_{(x,y)\to(0,0)} \frac{x^2-y^2}{x^2+y^2}.
\]

\item $\displaystyle \lim_{z\to 0}z^{-2}(1-\cos z))$ утгыг ол.

\item $f(z)=\frac{z^2}{|z|}$, $z\neq 0$, $f(0)=0$ гэж тодорхойлов.
$f$ функц 0 дээр тасралтгүй юу?

\item $\Ln z$ функц сөрөг бодит тэнхлэг дээр яагаад тасралтгүй биш вэ?
Тайлбарла.

\item Дараах хязгаар орших эсэхийг шалга:
\[
\lim_{z\to 0} \frac{xy}{x^2+y^2}.
\]

\item $f(z)=|z|$ функц $z_0=0$ дээр тасралтгүй эсэхийг $\varepsilon -\delta$ тодорхойлолтоор батал.

\item $\displaystyle \lim_{z\to 0} \frac{\sin z}{z}$ утгыг ол.
(Синусын тодорхойлолтыг ашиглаж бод.)

\item $\sqrt{z}$ функцийн бүх утгыг ол:
\[
z=-4.
\]
Мөн гол утгыг тодорхойл.

\end{enumerate}
\section*{Даалгаврын бодлого}

\begin{enumerate}

\item $\displaystyle \lim_{z\to 0} f(z)=L$ байх $\Longleftrightarrow $ нь $\forall t\in [0;2\pi[$ хувьд $\displaystyle \lim_{r\to 0} f(re^{it})=L$ байхыг батал. Үүнийг ерөнхий тохиолдолд  
$$\exists\displaystyle \lim_{z\to a} f(z) \Longleftrightarrow \displaystyle \lim_{r\to 0} f(re^{it})$$ хязгаар $t$-с үл хамааран оршин байх гэж бичиж болно. Энд $z-a=re^{it}.$ $\displaystyle\lim_{r\to 0} f(re^{it})=\lim_{r\to 0} u(r\cos t, r\sin t)+i\lim_{r\to 0} v(r\cos t, r\sin t)$ гэсэн нэг хувьсагчийн бодит функцийн хязгаар юм. 

\item 
\begin{equation}
\Arccos z= -i \Ln (z+\sqrt{z^2-1}) 
\end{equation} функцийн бодит ба хуурмаг хэсгийг ол. 

\item $\forall z\in\mathbb{C}$ хувьд $\cos(\Arccos(z))=z$ байхыг батал. Өөрөөр хэлбэл $\cos z$ урвуу функц олон утгат функц байна.

\item $-1\leq x\leq 1$ үед $\Arccos(x)$ нь бодит хувьсагчийн $\arccos(x)$ функц болохыг харуул.

\item $\sin z$ урвуу функц $\Arcsin(z)$ ол.

\item $\cos z, \sin z$ функцүүд $\mathbb{C}$ дээр тасралтгүй байхыг батал.

\item  $\tan z, \cot z$ функцүүд тасралтгүй байх олонлогийг ол.

\item $|z|\geq |Re(z)|$ батал.

\item $|z_1|+|z_2|\geq |z_1+z_2|\geq ||z_1|-|z_2||$ батал.

\item $|z|^2\geq 2Re(z)Im(z)$ батал.

\end{enumerate}
